%% =====================================================================
%%  T-18-02  Titles, Uniform and Trappings
%%  -------------------------------------------------------------------
%%  One idea per file. Core is mandatory. Slots in canonical order.
%%  Max 700 words. No layout commands. No filler. New source material
%%  for this entry goes in sources/<BOOK>/T-18-02.tex -- never by
%%  rewriting this file (docs/RULES.md R0.1).
%% =====================================================================
%% @kind: topic
%% @id: T-18-02
%% @title: Titles, Uniform and Trappings
%% @subtitle: The three costumes that are never checked
%% @chapter: c18
%% @order: 20
%% @status: stable
%% @sources: B4
%% @updated: 2026-09-19
%% @allow-rewrite: slot migration to the seven-question format (docs/RULES.md section 2)

\topic{T-18-02}{Titles, Uniform and Trappings}{The three costumes that are never checked}


\begin{core}
Three symbols do most of the work: a title, a uniform, and the trappings of position --- the car, the office, the cut of the clothes. All three are detachable from the thing they signify, all three are cheap, and none of them is routinely verified.
\end{core}

\begin{mechanism}
Each symbol answers a question that would otherwise take investigation, and each is answered in a form that discourages follow-up.

A title asserts rank and implies that the rank was conferred by someone competent to confer it. It is the most portable of the three because it requires no props --- only that it be said and not contradicted, and that it be specific enough to sound earned. Titles borrowed from adjacent fields work best, since the audience recognises the vocabulary without being able to check the credential.

A uniform asserts institutional backing and transfers the institution's authority to the wearer for the duration of the contact. It is the most powerful and the easiest to acquire, because the garment is widely available and the response to it is trained from childhood.

Trappings assert success, and success is read as evidence of competence and of honesty --- the assumption being that a person who has done well did not need to deceive. Trappings are rented rather than owned, which makes them the cheapest of the three to fake and the hardest to disprove at the moment of contact.
\end{mechanism}

\begin{application}
  \item Choose the title for its specificity and its adjacency to the field you need authority in.
  \item Wear the institutional marker where the contact is brief and verification is impractical.
  \item Rent the trappings for the occasion. The judgement is made in the first minutes and does not survive to be checked.
  \item Introduce the symbol early, then let the request arrive inside the frame it sets.
\end{application}

\begin{feedback}
  \item A title is offered but its conferring institution is not, or is named vaguely.
  \item The uniform or the setting appears only for this contact and not on any other occasion.
  \item The trappings are visible but the source of them is never mentioned.
  \item The symbol is introduced before any substantive claim.
  \item Questioning the symbol produces offence rather than information.
\end{feedback}

\begin{failure}
All three symbols are falsifiable and frequently falsified, so a person who relies on them is exposed the first time anyone checks. The uniform is the most dangerous of the three to misuse: wearing one you are not entitled to is a criminal offence in most jurisdictions, quite apart from whatever the contact itself amounts to. And the habit of verifying nothing transfers to the verifier --- someone who accepts symbols without checking becomes the reliable target that the technique is designed to find.
\end{failure}

\begin{countermeasures}
  \item Ask for the conferring body and verify it independently. A title that cannot name its source is a description, not a credential.
  \item Ask to see the credential rather than hearing about it --- an identifier, a registration, a callback to a number you found yourself.
  \item Never use contact details supplied by the person claiming authority. Find the institution independently and call it.
  \item Discount trappings entirely. Success displayed to you was selected to be displayed.
  \item Introduce a delay. Uniforms, titles and rented settings are all designed for a single contact and decay badly over two days.
\end{countermeasures}

\sources{B4}
\seesources
\seealso{T-18-01, T-23-03, T-22-02, T-09-01}
